% sec_introduction.tex -- Introduction
\section{Introduction}\label{sec:introduction}

The Cayley transform
\[
\iota(x)=\frac{1+\ii x}{1-\ii x}
\]
identifies the one-point compactification~$\overline\RR$ of the line
with the circle~$\TT$.  Under this identification normalized Haar
measure on~$\TT$ pulls back to the Cauchy law
\[
\frac{1}{\pi(1+x^{2})}\dd x=P_1(x)\dd x.
\]
This classical identity is the starting point of the paper.  It has a
geometric aspect, namely the uniqueness of the Cayley chart among
normalized orientation-preserving parametrizations of the circle, and it
has an analytic aspect, because it places the problem directly inside
the Poisson semigroup on~$\RR$.  The rigidity statement is proved in
Theorem~\ref{thm:haar-pullback-uniqueness}; the corresponding entropy
identity in Cayley coordinates is
Proposition~\ref{thm:precision-ledger}.

Our first theme concerns large-time entropy asymptotics for Poisson
averages.  If $\nu$ is a probability measure on~$\RR$ with mean
$\bar\gamma$, then after the rescaling $x=\bar\gamma+ty$ the ratio
between the shifted Poisson kernel and the reference Cauchy profile is
governed by
\[
R_\varepsilon(y):=\frac{1+y^2}{1+(y-\varepsilon)^2}.
\]
Writing $y=\tan\theta$ turns $R_\varepsilon$ into a finite trigonometric
object.  Theorem~\ref{thm:cayley-chebyshev-mode} makes this precise: the
Taylor coefficients of~$R_\varepsilon$ are given by an explicit
Cayley--Chebyshev mode calculus, with finite Fourier expansions in the
angle variable.  Once this structure is available, the parity of the
modes forces the asymptotic expansion of
\[
D_{\mathrm{KL}}(P_t*\nu\|P_t(\cdot-\bar\gamma))
\]
to involve only even powers of~$t^{-1}$; see
Theorem~\ref{thm:poisson-kl-even-orders}.

The even-order expansion can then be pushed to order~$t^{-8}$.
Theorem~\ref{thm:poisson-kl-eighth} gives the corresponding coefficient
formula in terms of the centered moments of~$\nu$.  The next step is to
recenter the coefficients at the symmetric two-point law with fixed
variance.  This produces an explicit defect hierarchy.
Theorem~\ref{thm:poisson-defect-ladder} rewrites the $t^{-6}$ and
$t^{-8}$ coefficients in terms of non-negative defect coordinates, while
Theorem~\ref{thm:poisson-defect-amplification} shows that the
eighth-order defect controls the sixth-order one with a strictly larger
coefficient, except at the symmetric two-point extremizer.  In
Theorem~\ref{thm:poisson-defect-stability} this is turned into a
quantitative $W_2$-stability statement.

The second theme is functional-analytic.  The secant family
\[
\Psi_\varepsilon(y):=\frac{R_\varepsilon(y)-1}{\varepsilon}
\]
has the exact Gram kernel
\[
K(a,b)=\frac{2}{4+(a-b)^2},
\]
as shown in Theorem~\ref{thm:mode-gram-kernel}.  Theorem~\ref{thm:mode-space-rkhs}
identifies the closed span of the secants with the zero-mean space
$L^2_0(\omega)$ and simultaneously with the reproducing-kernel Hilbert
space generated by~$K$.  Proposition~\ref{thm:gram-space-poisson-image}
then realizes this space as $\sqrt{\pi}\,e^{-|D|}L^2(\RR)$, so the
kernel space becomes a translation-invariant strip space with explicit
Fourier norm and holomorphic extension to
\[
\mathbb S=\{z\in\CC:\ |\Im z|<1\}.
\]
In Proposition~\ref{prop:poisson-hardy-splitting} the strip model splits
orthogonally into two Hardy components.

The lattice theory arises from the integer orbit of the kernel.
Theorem~\ref{thm:poisson-lattice-sampling} shows that the families
$\{K(\cdot,n)\}_{n\in\ZZ}$ and $\{\Psi_n\}_{n\in\ZZ}$ are Riesz bases
for their closed spans and that sampling on~$\ZZ$ is stable on the
lattice-generated subspace $\mathcal H_K(\ZZ)$.  Theorem~\ref{thm:poisson-cardinal-reconstruction}
then yields a cardinal kernel
$L\in\mathcal H_K(\ZZ)$ with $L(n)=\delta_{0n}$, a Shannon-type
reconstruction formula
\[
F_\alpha(z)=\sum_{n\in\ZZ}\alpha_nL(z-n),
\]
and an exact expression for the interpolation norm in terms of the
reciprocal lattice symbol.  In this way the strip realization leads to a
completely explicit principal shift-invariant model.

A further consequence is a centered reconstruction theorem from Poisson
data.  Theorem~\ref{thm:poisson-central-two-channel} shows that the
central Poisson trace and its central derivative recover the Laplace
transform of the centered characteristic function.  Laplace uniqueness
and characteristic-function inversion then recover the centered law, and
adding the mean recovers the original measure.  In the RKHS language
these two observation channels appear as evaluation functionals.

The basic analytic input is classical: the Poisson semigroup and its
Fourier representation
\cite{Folland1995,Grafakos2014,Stein1971}, Laplace-transform uniqueness
\cite{Widder1941}, characteristic-function inversion
\cite{Lukacs1970,Feller1971}, reproducing kernels
\cite{Aronszajn1950}, and Hardy-space boundary theory
\cite{Duren1970,Rudin1987}.  These ingredients are used here in a rather
concrete way, adapted to the Cayley normalization and the comparison
kernel~$R_\varepsilon$.

The entropy results should be viewed against the literature on the
entropic central limit theorem, entropy production, and stable-law
asymptotics; see
\cite{Barron1986,JohnsonBarron2004,ArtsteinBallBartheNaor2004,CarlenCarvalho1992}
and
\cite{BobkovChistyakovGotzeEdgeworth2013,BobkovChistyakovGotzeBerryEsseen2014,BobkovChistyakovGotze2013,Toscani2015,Toscani2016}.
The contribution of the present paper is not a new general
entropy-production inequality.  Rather, the Cayley normalization
identifies a Poisson/Cauchy model in which the coefficient calculus can
be carried out explicitly at each order.  In that model the parity of
the Cayley--Chebyshev modes forces the disappearance of odd terms, and
the order~$t^{-8}$ coefficient admits a defect decomposition strong
enough to detect the symmetric two-point law.

For the lattice part we rely on the standard theory of finitely
generated and principal shift-invariant spaces
\cite{deBoorDeVoreRon1994,RonShen1995,Bownik2000,AldroubiGrochenig2001},
together with later sampling refinements in reproducing-kernel and
totally positive settings
\cite{GrochenigRomeroStockler2018,GrochenigRomeroStockler2020,FuhrGrochenigHaimiKlotzRomero2018}.
Explicit interpolation formulae for Cauchy-type generators also appear
in \cite{KiselevMininNovikov2016,KiselevMininNovikov2019}, and recent
work on rational or closely related generators further emphasizes the
role of complex-analytic structure in sampling and Gabor theory
\cite{BelovKulikovLyubarskii2022,BelovKulikovLyubarskii2023,UlanovskiiZlotnikov2025}.
For a broader cardinal-interpolation perspective we also refer to
\cite{Buhmann2003}.  The present contribution is not a new abstract
sampling theorem; it is the identification of the Poisson-strip model
for the specific kernel~$K$ and the resulting closed formulas for the
lattice symbol, the cardinal kernel, and the interpolation norm.

\subsection{Theorem-by-theorem comparison with prior work}

Proposition~\ref{prop:cayley-homeo} is the standard boundary form of the
Cayley disk/half-plane correspondence from classical Hardy-space
theory \cite[Chapter~11]{Duren1970}.  Proposition~\ref{prop:phase-derivative}
is the usual tangent-to-Cauchy change of variables from probability
theory \cite[Chapter~II]{Feller1971}.  Theorem~\ref{thm:haar-pullback-uniqueness}
takes that classical pullback identity and adds a rigidity statement:
within the normalized orientation-preserving $C^1$ class, the pullback
condition forces the chart to be exactly~$\iota$.  Proposition~\ref{thm:precision-ledger}
is the standard differential-entropy change-of-variables formula in the
present notation \cite[Chapter~9]{CoverThomas2006}.  Proposition~\ref{prop:poisson-identities}
simply fixes the Poisson kernel normalization and Fourier multiplier
already classical in harmonic analysis \cite[Chapter~II]{Stein1971}.  Proposition~\ref{thm:poisson-semigroup-law}
is likewise supporting background on the Poisson semigroup and its
generator \cite[Chapter~II]{Stein1971}.

Proposition~\ref{thm:poisson-second-order} moves beyond the general
stable-law entropy literature \cite[Theorem~1.1]{BobkovChistyakovGotze2013}
by giving the explicit second-order comparison profile, the sharp
uniform ratio error, and the resulting $L^1$, $L^\infty$, and
Kullback--Leibler rates for the translated Poisson reference.  Theorem~\ref{thm:cayley-chebyshev-mode}
starts from the classical Chebyshev identities \cite[Chapters~1 and~5]{MasonHandscomb2002}
but identifies, for this specific comparison kernel, finite
Cayley--Chebyshev mode formulas with explicit parity and Fourier
expansions.  Proposition~\ref{prop:cayley-mode-remainder} is the
corresponding finite-truncation refinement, made explicit here as a
globally bounded rational remainder on~$\RR^2$.  Lemma~\ref{lem:mode-parity-vanish}
turns that mode parity into the exact weighted cancellation rule needed
for the entropy asymptotics.  Theorem~\ref{thm:poisson-kl-even-orders}
then sharpens the abstract expansion theory of entropic CLT papers
\cite[Theorem~1.1]{BobkovChistyakovGotzeEdgeworth2013} by proving that
every odd power disappears in the present Poisson/Cauchy model.  Theorem~\ref{thm:poisson-kl-eighth}
pushes that calculation to an exact order-$t^{-8}$ coefficient written
in centered moments.  Theorem~\ref{thm:poisson-defect-ladder} reorganizes
those coefficients into nonnegative defect coordinates adapted to the
symmetric two-point law.  Theorem~\ref{thm:poisson-defect-amplification}
adds the explicit coefficient comparison
$\Delta_8\ge (13\sigma^2/8)\Delta_6$ with equality only at the extremal
two-point law.  Theorem~\ref{thm:poisson-defect-stability} converts that
defect information into a quantitative $W_2$-stability estimate, which
is more explicit than the general entropy-production stability results
in \cite[Theorem~1]{CarlenCarvalho1992}.

Theorem~\ref{thm:mode-gram-kernel} takes the abstract RKHS viewpoint of
Aronszajn \cite[Theorem~1]{Aronszajn1950} and computes the exact Gram
kernel of the secant orbit as $2/(4+(\varepsilon-\eta)^2)$.  Corollary~\ref{cor:mode-quadratic-integrals}
extracts from that kernel the closed binomial formulas for the
quadratic mode integrals.  Theorem~\ref{thm:mode-space-rkhs} then proves
that the closed secant span is exactly $L^2_0(\omega)$ and not merely an
abstract RKHS model.  Theorem~\ref{thm:poisson-central-two-channel}
combines Laplace-transform uniqueness \cite[Chapter~VI]{Widder1941} with
characteristic-function inversion \cite[Chapter~2]{Lukacs1970} to show
that the concrete observation pair $(A,B)$ determines the centered law.
Corollary~\ref{cor:poisson-symmetrization} isolates the information
contained in $A$ alone, namely the centered symmetrization.  Corollary~\ref{cor:poisson-symmetric-case}
identifies symmetry exactly with the observable condition $B\equiv 0$.
Theorem~\ref{thm:poisson-channel-rkhs} places those two channels inside
the RKHS by identifying them as evaluations of the fluctuation profile
and of its image under the unitary map.  Proposition~\ref{thm:gram-space-poisson-image}
goes further than the abstract existence theorem by identifying this
particular RKHS with the explicit Poisson image space
$\sqrt{\pi}\,e^{-|D|}L^2(\RR)$ and its strip holomorphic extension.
Proposition~\ref{prop:poisson-hardy-splitting} is the standard
positive/negative frequency Hardy decomposition transported to this
normalization \cite[Chapter~11]{Duren1970}\cite[Chapter~17]{Rudin1987}.

Theorem~\ref{thm:poisson-lattice-sampling} specializes the general
principal shift-invariant theory of Ron and Shen
\cite[Theorem~2.2.14]{RonShen1995} by computing the exact lattice symbol
for the Cauchy kernel and hence the explicit Riesz bounds
$A_{\ZZ},B_{\ZZ}$.  Theorem~\ref{thm:poisson-cardinal-reconstruction}
then strengthens the standard interpolation formula
\cite[formula~(2.2.11)]{RonShen1995} by giving the exact Fourier formula
for the cardinal kernel and the exact RKHS norm of every interpolant.

The paper is organized as follows.  Section~\ref{sec:preliminaries}
fixes notation and conventions.  Section~\ref{sec:cayley-gate} records
the Cayley parametrization, and Section~\ref{sec:haar-pullback} proves
the uniqueness of the Haar pullback.  Section~\ref{sec:precision-potential}
develops the entropy identity, the mode calculus, the even-order
expansion, and the defect estimates.  Section~\ref{sec:gram-space}
constructs the Gram space and proves the centered reconstruction result.
Section~\ref{sec:strip-lattice} gives the strip realization, the Hardy
splitting, and the lattice sampling theory.  The appendix contains the
weighted trigonometric integrals needed for the explicit eighth-order
coefficient.
